
Sebagai penutup, tugas kecil 2 ini menantang kami untuk menyelesaikan permasalahan geometri komputasional secara terstruktur, mulai dari membaca model 3D berformat OBJ, membangun representasi ruang menggunakan octree, hingga melakukan proses voxelisasi dan visualisasi hasilnya. Berbeda dengan pendekatan yang sepenuhnya naif (Tucil 1 Algoritma Brute Force), pada tugas ini kami belajar bahwa pemilihan strategi dan struktur data yang tepat sangat menentukan efisiensi, terutama ketika ruang pencarian semakin besar seiring peningkatan nilai max-depth. Proses implementasi, pengujian pada berbagai kasus dasar maupun edge case, serta analisis hasil membuat kami lebih memahami trade-off antara ketelitian hasil dan waktu komputasi. Menurut kami, pengalaman ini merupakan langkah penting untuk melatih pola pikir algoritmik, ketelitian dalam rekayasa perangkat lunak, dan kemampuan evaluasi solusi secara menyeluruh. Harapannya, melalui tugas-tugas berikutnya, kemampuan kami dalam merancang solusi yang benar, efisien, dan robust dapat semakin matang. \textbf{Tugas ini disusun sepenuhnya tanpa bantuan kecerdasan buatan (Generative AI), melainkan hasil pemikiran dan analisis mandiri.}

\vspace{10px}
\begin{flushright}
    Bandung, 27 Maret 2026
    \vspace{3px}

    
\includegraphics[width=0.4\textwidth]{sign-aufar.png}

    \vspace{3px}
    Muhammad Aufar Rizqi Kusuma\\
    13524061

    \includegraphics[width=0.2\textwidth]{sign-athilla.png}

    \vspace{3px}
    Athilla Zaidan Zidna Fann\\
    13524068
\end{flushright}